% TOP_PROBLEM: {"id": 3340, "title": "Blatter-Specker Theorem for ternary relations", "category_id": 18, "category": {"id": 18, "name": "logic", "display_name": "Logic", "description": "Problems in mathematical logic, model theory, proof theory, and finite model theory.", "slug": "logic", "order_index": 18, "created_at": "2026-07-31T15:26:25.671Z"}, "status": "open"}
% FIRST_POSTED: null
% ENTRY_KIND: statement_only
% Imported from: batch03.tex; UnsolvedMath ID 3340
\documentclass[11pt]{article}
\usepackage[margin=25mm]{geometry}
\usepackage{fontspec}
\setmainfont{Latin Modern Roman}
\usepackage{amsmath,amssymb,mathtools}
\usepackage{unicode-math}
\setmathfont{Latin Modern Math}
\usepackage{xurl,hyperref}
\hypersetup{colorlinks=true,urlcolor=blue}
\setcounter{secnumdepth}{0}
\setlength{\parindent}{0pt}
\setlength{\parskip}{5pt}
\setlength{\emergencystretch}{3em}
\newcommand{\sref}[2]{\hyperlink{source:#1:#2}{[S#2]}}
\newcommand{\cataloguescope}[1]{\paragraph{Recorded target.}#1}
\begin{document}
% UnsolvedMath ID 3340; source catalogue code OPG-37444
\setcounter{section}{3339}
\section{Blatter-Specker Theorem for ternary relations}\label{3340}
{\small\sffamily UnsolvedMath ID 3340 · Logic}
\subsection{Definitions and mathematical statement}

\paragraph{Shared notation.}$\mathbb N=\{1,2,\ldots\}$, and $\mathbb Z,\mathbb Q,\mathbb R,\mathbb C$ denote the integers, rationals, reals and complex numbers. Any other conventions are specified locally.


Let $\tau$ be a fixed finite relational vocabulary with each relation symbol of arity at most three, and let $\varphi$ be a monadic second-order (MSO) sentence over $\tau$; MSO quantifies over elements and sets of elements. Let $C$ be its class of finite models. For $n\geq0$, define $f_C(n)$ as the number of $\tau$-structures on the labeled domain $\{0,\ldots,n-1\}$ satisfying $\varphi$. Ultimate periodicity modulo $m\geq1$ means that there exist integers $N\geq0$ and $p\geq1$ such that $f_C(n+p)\equiv f_C(n)\pmod m$ for every $n\geq N$.
\paragraph{Statement recorded in the supplied source.}
Does every such MSO-definable class $C$ have $f_C$ ultimately periodic modulo every positive integer $m$, when ternary relation symbols are allowed? This asks whether the Specker--Blatter theorem for unary and binary vocabularies extends to arity three. \sref{3340}{2}
\subsection{Short English statement}
Does the eventual modular periodicity of labeled counts for monadic second-order definable classes survive the addition of ternary relations?
\subsection{Sources}
\begin{description}
\item[\hypertarget{source:3340:1}{S1}] ULAM.AI and the UnsolvedMath Contributors, UnsolvedMath: A Curated Collection of Open Mathematics Problems, record 3340 (OPG-37444). CC BY 4.0. \url{https://huggingface.co/datasets/ulamai/UnsolvedMath}.
\item[\hypertarget{source:3340:2}{S2}] Open Problem Garden, Blatter-Specker Theorem for ternary relations, node 37444, original problem page. Original question and full discussion. \url{https://www.openproblemgarden.org/op/blatter_specker_theorem_for_ternary_relations}.

\end{description}
\label{end:3340}
\end{document}
